% from ega4-1-fr.tex
\oldpage[IV-1]{247}

\begin{proof}
On sait que la condition est nécessaire ((1.8.2) et (1.9.5, (vi)))~; pour
montrer qu'elle est suffisante, on est ramené au cas où $X$ est affine, en
vertu de (1.9.5, (iv))~; en outre, en vertu de (1.9.11), on peut se borner au
cas où $f^{-1}(E)$ est pro-constructible, \emph{ou} au cas où $f^{-1}(E)$ est
ind-constructible. Si $f$ est surjectif et quasi-compact, et $f^{-1}(E)$
pro-constructible, il résulte de (1.9.5, (vii)) que
$E=f(f^{-1}(E))$ est pro-constructible. Si $f$ est surjectif et localement de
présentation finie, et $f^{-1}(E)$ ind-constructible,
$E=f(f^{-1}(E))$ est ind-constructible par (1.9.5, (viii)), ce qui achève la
démonstration.
\end{proof}

\begin{env}[1.9.13]
\label{IV.1.9.13-fr}
Soient $X$ un préschéma, $\mathscr T$ sa topologie~; il résulte de (1.9.5,
(i) et (ii)) que les parties ind-constructibles (resp. pro-constructibles) de
$X$ sont les parties ouvertes (resp. fermées) pour une topologie sur $X$,
appelée \emph{topologie constructible} et que nous noterons
$\mathscr T^{\mathrm{cons}}$~; nous désignerons par $X^{\mathrm{cons}}$
l'ensemble $X$ muni de la topologie $\mathscr T^{\mathrm{cons}}$.
\end{env}

\begin{proposition}[1.9.14]
\label{IV.1.9.14-fr}
Soient $X$ un préschéma, $\mathscr T$ sa topologie,
$\mathscr T^{\mathrm{cons}}$ la topologie constructible sur $X$.
\begin{enumerate}
  \item[(i)] La topologie $\mathscr T^{\mathrm{cons}}$ est plus fine que
  $\mathscr T$.
  \item[(ii)] Les parties localement constructibles de $X$ sont identiques
  aux parties à la fois ouvertes et fermées de l'espace $X^{\mathrm{cons}}$.
  \item[(iii)] Pour tout morphisme $f:X\to Y$, l'application sous-jacente de
  $X^{\mathrm{cons}}$ dans $Y^{\mathrm{cons}}$ est continue~; on la note
  $f^{\mathrm{cons}}$.
  \item[(iv)] Si le morphisme $f:X\to Y$ est quasi-compact,
  $f^{\mathrm{cons}}$ est une application fermée~; en particulier, si $f$
  est quasi-compact et bijectif, $f^{\mathrm{cons}}$ est un homéomorphisme.
  \item[(v)] Si un morphisme $f:X\to Y$ est localement de présentation
  finie, $f^{\mathrm{cons}}$ est une application ouverte~; en particulier,
  si $f$ est bijectif et localement de présentation finie,
  $f^{\mathrm{cons}}$ est un homéomorphisme.
  \item[(vi)] Pour tout ouvert $U$ de $X$, la topologie induite par
  $\mathscr T^{\mathrm{cons}}$ sur $U$ est identique à la topologie de
  $U^{\mathrm{cons}}$.
\end{enumerate}
\end{proposition}

\begin{proof}
En effet, (i) résulte de ce que tout ouvert pour $\mathscr T$ est un ouvert
pour $\mathscr T^{\mathrm{cons}}$ (1.9.6), et (ii) est une traduction de
(1.9.11)~; (iii), (iv) et (v) traduisent respectivement (1.9.5), (vi), (vii)
et (viii). Enfin pour démontrer (vi), il suffit de remarquer que l'injection
canonique $j:U\to X$ est localement de présentation finie (1.4.3), donc
$j^{\mathrm{cons}}:U^{\mathrm{cons}}\to X^{\mathrm{cons}}$ est une
application continue ouverte.
\end{proof}

\begin{proposition}[1.9.15]
\label{IV.1.9.15-fr}
Soit $X$ un préschéma.
\begin{enumerate}
  \item[(i)] Pour que $X^{\mathrm{cons}}$ soit quasi-compact, il faut et il
  suffit que $X$ soit quasi-compact.
  \item[(ii)] Pour que $X^{\mathrm{cons}}$ soit séparé, il faut et il suffit
  que $X$ soit quasi-séparé, et alors $X^{\mathrm{cons}}$ est localement
  compact et totalement discontinu.
  \item[(iii)] Pour que $X^{\mathrm{cons}}$ soit compact, il faut et il suffit
  que $X$ soit quasi-compact et quasi-séparé.
  \item[(iv)] Tout point de $X$ admet, pour la topologie
  $\mathscr T^{\mathrm{cons}}$, un voisinage ouvert et compact.
  \item[(v)] Pour qu'un morphisme $f:X\to Y$ soit quasi-compact, il faut et il
  suffit que l'application continue
  $f^{\mathrm{cons}}:X^{\mathrm{cons}}\to Y^{\mathrm{cons}}$ soit propre.
\end{enumerate}
\end{proposition}

\begin{proof}
(i) Comme $\mathscr T^{\mathrm{cons}}$ est plus fine que $\mathscr T$, il est
clair que si $X^{\mathrm{cons}}$ est quasi-compact, il en est de même de $X$~;
la réciproque résulte de (1.9.9).

(ii) Supposons $X$ quasi-séparé, et montrons que $X^{\mathrm{cons}}$ est
totalement discontinu~: en effet, si $x$, $y$ sont deux points distincts de
$X$ et si par exemple $x\notin\overline{\{y\}}$, il existe un
